% ============================================================
%  PipedPeer -- conference paper
%  Sardar Patel Institute of Technology
% ============================================================
\documentclass[conference]{IEEEtran}
\IEEEoverridecommandlockouts

\usepackage{cite}
\usepackage{amsmath,amssymb,amsfonts}
\usepackage{algorithmic}
\usepackage{graphicx}
\usepackage{textcomp}
\usepackage{xcolor}
\usepackage{booktabs}
\usepackage{tabularx}
\usepackage{url}
\usepackage{caption}
\usepackage{subcaption}

\newcolumntype{L}[1]{>{\raggedright\arraybackslash}p{#1}}
\newcolumntype{C}[1]{>{\centering\arraybackslash}p{#1}}
\newcommand{\figpath}[1]{../diagrams/#1}

\begin{document}

\title{PipedPeer: Zero-Configuration Distributed Computing via Recursive, Decentralised Coordination}

\author{
\IEEEauthorblockN{Yateen Vaviya}
\IEEEauthorblockA{\textit{Department of Computer Engineering} \\
\textit{Sardar Patel Institute of Technology}\\
Mumbai, Maharashtra, India \\
yateen.vaviya23@spit.ac.in}
\and
\IEEEauthorblockN{Vinayak Yadav}
\IEEEauthorblockA{\textit{Department of Computer Engineering} \\
\textit{Sardar Patel Institute of Technology}\\
Mumbai, Maharashtra, India \\
vinayak.yadav23@spit.ac.in}
\and
\IEEEauthorblockN{Sahil Gupta}
\IEEEauthorblockA{\textit{Department of Computer Engineering} \\
\textit{Sardar Patel Institute of Technology}\\
Mumbai, Maharashtra, India \\
sahil.gupta23@spit.ac.in}
}

\maketitle

\begin{abstract}
Distributed computing frameworks require a cluster that somebody has already
configured, and they require the program to be rewritten against their API. Both
costs are paid before any speedup is obtained, which puts the idle machines on a
small team's network out of reach for ordinary work. We present PipedPeer, a
system that runs an unmodified Python script on another machine on the network
with no cluster configuration and nothing installed on any participant beyond a
single binary. Two mechanisms make this possible. First, the script's
environment is built as a content-addressed Nix closure against a pinned package
set, so the same script yields an identical store path on every machine and on
any day; that path is then used as a cluster-wide cache key, turning environment
transfer from a per-job cost into a once-per-cluster cost. On a three-node
cluster the first submission of a job took 12.00\,s and subsequent submissions
of the same environment took 0.65\,s, an eighteenfold reduction. Second, the
parallel primitives the script already uses are substituted at import time by a
shim injected into the job's interpreter, and are distributed only when a cost
model fed by a measured bandwidth probe finds that moving the data costs less
than computing it locally. The model declines to distribute dense matrix
multiplication, and the project treats that refusal as a tested property rather
than a missing feature. Every remote path degrades to a working local one: a
worker killed during a distributed parallel map left the run to complete with
correct results.
\end{abstract}

\begin{IEEEkeywords}
distributed computing, peer-to-peer, transparent interception, reproducible
environments, content-addressed caching, cost-based scheduling, fault tolerance
\end{IEEEkeywords}

\section{Introduction}

The gap between running \texttt{python script.py} and running the same work on
several machines is large in operational effort and often small in the
computation involved. Apache Spark~\cite{zaharia2010spark},
Ray~\cite{moritz2018ray} and Dask~\cite{rocklin2015dask} all presuppose a
configured cluster with a head node, matching software on every worker, and a
scheduler process. Remote execution over \texttt{ssh} avoids that setup but
provides no environment management, no placement, no isolation and no fault
tolerance.

Two distinct obstacles are usually treated separately. The \emph{environment}
obstacle is that a remote machine must be made to match the local one, which is
the bulk of the work and the reason ad hoc remote execution does not scale past
one machine. The \emph{code} obstacle is that using more than one machine
normally requires rewriting the program.

PipedPeer removes both. Its contributions are:

\begin{enumerate}
\item \textbf{Peer-to-peer closure distribution keyed by store path.}
Content-addressed store paths and substitution are Nix's, not ours
\cite{dolstra2004nix}; what we add is pushing a closure between peers rather
than pulling from a cache, a two-predicate presence check that is correct both
for per-node stores and for a shared store, and the use of that check to decide
which peers are eligible to receive distributed work (Section~\ref{sec:env}).
\item \textbf{Transparent interception under a measured cost model.} Parallel
primitives are substituted at import time and distributed only when the
arithmetic favours it, with link bandwidth measured rather than assumed
(Sections~\ref{sec:intercept}, \ref{sec:cost}).
\item \textbf{A shuffle-only aggregation path.} Hash partitioning on the key
is standard \cite{dean2004mapreduce}; the choice here is to use \emph{only}
that path and never a combiner, which keeps every aggregation exact in value,
including user-defined ones, at the cost of always paying a full shuffle
(Section~\ref{sec:shuffle}).
\item \textbf{Placement without a scheduler.} Hard filters followed by a bounded
score, evaluated in the submitting client, with a three-phase lease providing
atomic reservation and recovery through expiry (Section~\ref{sec:placement}).
\item \textbf{A correctness-preserving fallback on every distributed path},
structural for the parallel-map race and by exception handling elsewhere, so a
failed or absent cluster changes the running time but never the result
(Section~\ref{sec:ft}).
\end{enumerate}

\section{Related Work}

Spark~\cite{zaharia2010spark}, Ray~\cite{moritz2018ray} and
Dask~\cite{rocklin2015dask} are mature and considerably more capable at scale
than the system described here; they differ in requiring a configured cluster
and, for Spark and to a lesser degree the others, a program written against
their API.

Two of them address the same two problems we do, and the distinction is narrow
enough to state precisely. Ray ships a drop-in
\texttt{ray.util.multiprocessing.Pool} and a \texttt{joblib} backend, and its
runtime environments ship and cache per-job dependencies; Modin~\cite{petersohn2020modin}
provides distributed pandas behind the pandas API. Each requires the user to
change an import line and to have a cluster already running. What is different
here is that no line of the user's program changes at all, because substitution
happens in \texttt{sitecustomize} before the first user statement executes, and
that the decision to distribute is gated on a measured cost model rather than
taken whenever the API is used. We are not aware of prior work that combines
import-time substitution of unmodified code with a per-call cost gate. BOINC~\cite{anderson2004boinc} pioneered volunteer computing but is a
platform for specific projects, each with its own server infrastructure, rather
than a general executor. Nix~\cite{dolstra2004nix} contributes the
content-addressed store on which our caching argument rests.

Work on offloading in edge computing addresses the same decision our cost model
makes. Deng et al.~\cite{deng2023} formulate offloading as a comparison of
transmission against local computation time; Yong et al.~\cite{yong2023} show
partial offloading to beat all-or-nothing strategies, which is the regime our
interception operates in. Learned policies~\cite{he2024} exceed what a
deterministic rule achieves, at the cost of a training phase that a
zero-configuration tool cannot assume. Zhao et al.~\cite{zhao2023} show
topology-aware placement to outperform scalar scoring, a limitation our
placement function shares and which we identify as future work.

On failure transparency, Bergstr\"om~\cite{kth2025thesis} surveys abstractions
that hide distributed failure from application code; since a user script here
contains no failure handling at all, this is the property Section~\ref{sec:ft}
must provide. The Python Array API~\cite{arrayapi2023} shares our goal of
portability without code change but inverts the mechanism, asking libraries to
converge on one interface rather than substituting the implementations behind
interfaces users already call.

\section{System Design}

Every machine runs one daemon; the command-line client is short-lived. The
coordinator is a library linked into that client rather than a service, so there
is no scheduler process, no head node and no master. A machine is a submitter
and a worker simultaneously, and a new participant joins by running the binary.

\begin{figure}[t]
\centering
\includegraphics[width=\columnwidth]{\figpath{fig_3.1_system_design.png}}
\caption{System design. The submitting side builds the environment and places
the job; the worker admits it, caches the closure and runs it sandboxed. Both
roles are the same binary.}
\label{fig:arch}
\end{figure}

\subsection{Placement}
\label{sec:placement}

Candidates are gathered from a local store of known peers, from UDP probes on
the subnet, and optionally from a registry. Five hard filters run in order:
architecture must match the closure, a GPU must exist if required, and free
memory, cores and VRAM must each satisfy the request. Survivors are scored:

\begin{equation}
\label{eq:score}
\begin{split}
S = \Bigl(1 &- \tfrac{u_{\text{cpu}}}{200} - \tfrac{u_{\text{mem}}}{200}
     - 0.05\,n_{\text{jobs}} \\
    &+ 0.10(\hat{c}-\tfrac12) + 0.06(\hat{f}-\tfrac12) \\
    &+ 0.10(\hat{m}-\tfrac12) + 0.10(\hat{r}-\tfrac12)\Bigr)\cdot H
\end{split}
\end{equation}

where $u$ are utilisations in percent, $n_{\text{jobs}}$ is the running job
count, $\hat{c},\hat{f},\hat{m},\hat{r}$ are core count, clock, free memory and
free-core ratio normalised onto $[0,1]$, and $H$ is a health score. Two
properties matter more than the coefficients. Load dominates hardware: the load
terms can subtract $1.0$ while the hardware terms move the score by at most
$0.36$, so a fast but busy machine loses to a slower idle one. And a missing
metric normalises to $0.5$, the neutral midpoint, so uneven telemetry across
heterogeneous machines neither rewards nor penalises a node.

Work is reserved by a three-phase lease. The admission check and the lease
insertion occur under one lock, so concurrent submitters cannot both claim the
last slot, and the concurrency limit counts reserved and running leases together
so a burst cannot slip through the window between the phases. Recovery needs no
further mechanism: a submitter that dies stops renewing and its slot is
reclaimed, while a worker that dies drops out of the healthy set and its job is
rescheduled.

\subsection{Environment Reproduction}
\label{sec:env}

The client scans the script's imports, resolves them to packages, generates a
Nix flake and builds a closure. The flake pins an \emph{exact} package-set
revision rather than a branch, and this is the load-bearing decision of the
entire design. A branch resolves at build time, so two machines building the
same script on different days would produce different store paths; because the
cache is keyed on the store path, the system would re-ship a large archive for
every job. An exact pin makes the store path a stable, content-derived name for
an environment across every machine and across time.

Before uploading, the client asks the target whether it already holds that path.
Two distinct predicates are needed and both are implemented: whether the archive
is cached, and whether the closure is materialised on disk. A deployment with a
shared store volume satisfies the second without the first, and a per-node store
satisfies the first before the second; either predicate alone is wrong for one
of the two layouts. After an upload the worker pushes the closure to peers that
lack it, which is what makes those peers eligible to receive distributed work.

\subsection{Transparent Interception}
\label{sec:intercept}

A Python module embedded in the binary is appended to the job's workspace
archive and placed on the interpreter's module path, so it is imported
automatically before the user's first line. Nothing is installed on any node and
the user's project on disk is never modified. Each library is patched on its
first import rather than up front, since importing a large framework such as
torch costs on the order of a second and charging that to jobs that never use it
would violate the guarantee of Section~\ref{sec:ft}.

Substituted primitives are parallel maps, pandas grouping, joining and chunked
reading, NumPy matrix multiplication and decompositions, \texttt{joblib} and
scikit-learn parallelism, and data-parallel training.

\subsection{Exact Aggregation Without Combiners}
\label{sec:shuffle}

For grouping and joining, rows are bucketed by a hash of the key. Equal keys
hash equally, so every key lands complete on exactly one node. Each node's
aggregation is therefore a complete aggregation over the keys it holds, and the
origin combines partials by concatenation alone. No combiner algebra is
involved, and consequently \emph{any} aggregation specification is exact in
value, including user-defined reductions that admit no combiner. Systems that
push combiners to the workers also fall back to a full shuffle for
non-combinable aggregations, so this is not a stronger guarantee than theirs; it
is the same guarantee reached by always taking the slower path, which suits a
shim that cannot inspect the user's aggregation function. Exactness here is of
values, not row order: concatenating buckets does not reproduce pandas'
documented join ordering.

\subsection{Data-Parallel Training Without a Socket Mesh}

Standard data-parallel training establishes direct connections between every
pair of ranks on ports of its own choosing, an assumption that holds in a data
centre and often nowhere else. Instead, each rank posts its gradient to the lead
node's ordinary daemon port and receives the full set once every rank has
arrived; the daemon never interprets the payload and averaging happens in the
shim. This trades bandwidth for requiring nothing beyond the port the rest of
the system already uses.

\section{The Cost Model}
\label{sec:cost}

Interception is only worthwhile when the cluster wins, so every intercepted call
is gated. Writing $b$ for payload bytes, $\phi$ for arithmetic intensity in
operations per byte, $B$ for measured bandwidth, $R$ for local throughput and
$K$ for node count, work is distributed only when

\begin{equation}
\label{eq:cost}
\frac{b\phi}{R} \;>\; \frac{b}{B} \;+\; \frac{b\phi}{RK}\times 1.3
\end{equation}

the factor $1.3$ charging 30\% overhead against ideal parallel speedup. Three
guards precede the arithmetic: at least two nodes, a payload of at least 32\,MB,
and a carve-out rejecting low-intensity work below 512\,MB. Bandwidth is
measured by pushing 64\,MB through the real dispatch path, cached for five
minutes; if the probe fails the answer is to stay local.

\begin{figure}[t]
\centering
\includegraphics[width=\columnwidth]{\figpath{fig_4.9_cost_model.png}}
\caption{Decision regions of Equation~\ref{eq:cost} at three nodes on a
500\,MB/s link.}
\label{fig:cost}
\end{figure}

A separately calibrated model governs dense linear algebra and produces a result
worth emphasising: \textbf{matrix multiplication is never distributed}. Local
BLAS sustains roughly 200\,GFLOP/s, which no cluster transport beats at
practical sizes. The project encodes this as a tested property, failing its
demonstration suite if a matrix multiplication is ever observed leaving a node.
Singular value decomposition, two orders of magnitude slower locally, does move.

\section{Degradation Guarantee}
\label{sec:ft}

Using the cluster must never change the result, and must not leave the caller
worse off than a local run when the cluster is slow or absent. The timing half
of that is bounded rather than absolute: Section~\ref{sec:res-shim} measures the
cost of leaving interception on at a few per cent of a short job, and only the
parallel-map path has a structural bound. Five layers apply. Chunks of a parallel map are dealt alternately to the local side
and the cluster and fill a result table first-wins, with idle local cores
re-running any chunk the cluster has not returned; the call returns only once
the local side has covered every item, so a wholly unresponsive cluster degrades
to a plain local pool. Beneath that, every intercepted call falls back to the
original implementation on any exception, a failed peer is skipped in favour of
the next, a dead worker process is respawned, and a failed job is rescheduled
with exponential backoff.

\section{Evaluation}

\subsection{Method and its limits}

Measurements were taken on one host with 20 logical cores running three worker
daemons in containers. \textbf{This bed cannot measure multi-machine speedup},
because moving processor-bound work between containers on one host cannot make
that host faster, and no speedup figure is claimed. It measures correctness,
overhead, cache behaviour and fault tolerance, which are the properties the
design argument rests on.

\subsection{Environment cache}

The same job submitted five times took 12.00\,s on the first submission and a
median of 0.65\,s thereafter, an \textbf{eighteenfold reduction}
(Fig.~\ref{fig:cache}). The first submission ships the closure; the rest find it
materialised and ship only the project files. This is the direct consequence of
pinning the package set so that a store path names an environment identically
everywhere.

\begin{figure}[t]
\centering
\includegraphics[width=\columnwidth]{\figpath{fig_5.1_closure_cache.png}}
\caption{Cold and warm submission of the same environment.}
\label{fig:cache}
\end{figure}

\subsection{Cost of leaving interception on}
\label{sec:res-shim}

Since interception is on by default, its cost when the model declines to
distribute is what most users pay. On a workload the model declines, a whole
remote job took 0.93\,s against 0.85\,s with interception disabled, an overhead
of \textbf{9\%}.

Interpreter startup for a job that imports nothing heavy is a separate matter.
Eagerly patching libraries cost $21.60\times$ the unshimmed startup, against a
budget of $3\times$ the project sets itself; deferring each library's patching
to its first import brings this to $1.20\times$.

\subsection{Correctness and fault tolerance}

The suite passes 246 test cases with none failing. Distribution correctness is
established by comparison against local execution: grouped aggregation is
frame-equal to local pandas, the four join types are equal up to row order (a
concatenation of shuffle buckets does not reproduce pandas' join ordering), and
each rank's gradient in data-parallel training is exactly the mean across
ranks.

Under fault injection, a twenty-item parallel map was distributed across the
cluster and a worker was killed mid-computation. The run completed with correct
results, the checksum matching the serial computation.

\begin{table}[t]
\centering
\caption{PipedPeer against existing systems. The lower block is where it is
the weaker system.}
\label{tab:cmp}
\renewcommand{\arraystretch}{1.2}
\footnotesize
\begin{tabularx}{\columnwidth}{X C{0.75cm} C{0.7cm} C{0.6cm} C{0.7cm} C{0.7cm}}
\toprule
\textbf{Property} & \textbf{Ours} & \textbf{Spark} & \textbf{Ray} & \textbf{Dask} & \textbf{ssh} \\
\midrule
\multicolumn{6}{l}{\textit{What it does differently}} \\
Runs unmodified Python     & yes & no  & part. & part. & yes \\
No cluster to configure    & yes & no  & no    & no    & part. \\
No head or master node     & yes & no  & no    & no    & yes \\
Reproduces the environment & yes & part. & part. & part. & no \\
Sandboxes each job         & part. & no & no   & no    & no \\
Survives losing a node     & yes & yes & yes   & yes   & no \\
\midrule
\multicolumn{6}{l}{\textit{Where it is behind}} \\
Authenticated, encrypted   & \textbf{no} & yes & part. & part. & yes \\
Works beyond one LAN       & part. & yes & yes & yes & yes \\
Runtimes beyond Python     & \textbf{no} & yes & part. & no & yes \\
Proven in production       & \textbf{no} & yes & yes & yes & yes \\
\bottomrule
\end{tabularx}
\end{table}

The upper block lists the criteria the system was designed to change, so it
wins them by construction and they should be discounted accordingly. The lower
block is the informative half. Ours is the only system in the table with no
authentication and no transport encryption whatsoever, which confines it to an
already-trusted network; it is also the only one limited to one language and one
subnet, and the only one without a production track record. Spark, Ray and Dask
are substantially more capable at scale. The claim is not that this system is
better, but that it occupies a different point on the curve, trading reach,
security and maturity for the removal of setup cost and code change.

\subsection{Threats to validity}

The single-host bed is the principal threat and the reason no speedup is
claimed. Container workers share a processor, a memory system and the loopback
interface, so both transfer costs and contention differ from a real network. The
cache result is the least affected, since it turns on whether an archive is sent
at all rather than on how fast it travels; the interception overhead figures are
the most affected and should be read as lower bounds on the benefit and upper
bounds on the relative overhead.

\section{Conclusion}

PipedPeer runs unmodified Python on other machines with no cluster to configure,
by reproducing environments as content-addressed closures whose store path
serves as a cluster-wide cache key, and by substituting the parallel primitives
a script already uses under a cost model driven by measured bandwidth. The
measurements that carry the argument are an eighteenfold reduction in repeat
submission cost and correct completion of a distributed computation after a
worker was killed mid-run.

The part we consider most transferable is the discipline of the cost model. A
system that distributed everything it could intercept would be slower than plain
Python on most real workloads, because a modern BLAS or a warm page cache
usually beats the network. Measuring rather than assuming, declining to
distribute when local wins, and treating that refusal as a property to be tested
is what makes always-on interception defensible.

Future work is led by authentication and an encrypted transport, without which
the system is confined to a trusted network; then discovery beyond the subnet,
locality- and topology-aware placement, and measurement across separate machines
to obtain the speedup figures this evaluation deliberately does not claim.

\begin{thebibliography}{00}

\bibitem{zaharia2010spark} M.~Zaharia, M.~Chowdhury, M.~J.~Franklin, S.~Shenker, and I.~Stoica, ``Spark: Cluster computing with working sets,'' in \textit{Proc. 2nd USENIX Workshop on Hot Topics in Cloud Computing}, 2010.

\bibitem{moritz2018ray} P.~Moritz et al., ``Ray: A distributed framework for emerging AI applications,'' in \textit{Proc. 13th USENIX OSDI}, 2018, pp.~561--577.

\bibitem{rocklin2015dask} M.~Rocklin, ``Dask: Parallel computation with blocked algorithms and task scheduling,'' in \textit{Proc. 14th Python in Science Conf.}, 2015, pp.~130--136.

\bibitem{anderson2004boinc} D.~P.~Anderson, ``BOINC: A system for public-resource computing and storage,'' in \textit{Proc. 5th IEEE/ACM Workshop on Grid Computing}, 2004, pp.~4--10.

\bibitem{dolstra2004nix} E.~Dolstra, M.~de~Jonge, and E.~Visser, ``Nix: A safe and policy-free system for software deployment,'' in \textit{Proc. 18th LISA}, 2004, pp.~79--92.

\bibitem{deng2023} Y.~Deng, Z.~Chen, X.~Chen, and Y.~Fang, ``Task offloading in multi-hop relay-aided multi-access edge computing,'' \textit{IEEE Trans. Veh. Technol.}, vol.~72, no.~1, pp.~1372--1376, 2023.

\bibitem{yong2023} D.~Yong, R.~Liu, X.~Jia, and Y.~Gu, ``Joint optimization of multi-user partial offloading strategy and resource allocation in D2D-enabled MEC,'' \textit{Sensors}, vol.~23, no.~5, 2023.

\bibitem{he2024} H.~He, X.~Yang, X.~Mi, H.~Shen, and X.~Liao, ``Multi-agent deep reinforcement learning based dynamic task offloading in a D2D mobile-edge computing network,'' \textit{Sensors}, vol.~24, no.~16, 2024.

\bibitem{zhao2023} Z.~Zhao, J.~Perazzone, G.~Verma, and S.~Segarra, ``Congestion-aware distributed task offloading in wireless multi-hop networks using graph neural networks,'' arXiv:2312.02471, 2023.

\bibitem{petersohn2020modin} D.~Petersohn et al., ``Towards scalable dataframe systems,'' \textit{Proc. VLDB Endowment}, vol.~13, no.~12, pp.~2033--2046, 2020.

\bibitem{dean2004mapreduce} J.~Dean and S.~Ghemawat, ``MapReduce: Simplified data processing on large clusters,'' in \textit{Proc. 6th USENIX OSDI}, 2004, pp.~137--150.

\bibitem{kth2025thesis} A.~Bergstr\"om, ``Programming models for failure-transparent distributed systems,'' M.S. thesis, KTH Royal Institute of Technology, 2025.

\bibitem{arrayapi2023} A.~Meurer et al., ``Python array API standard: Toward array interoperability in the scientific Python ecosystem,'' in \textit{Proc. SciPy}, 2023.

\end{thebibliography}

\end{document}
